\section{Amplification and the forcing consequence}
\label{sec:consequences}

Theorem~\ref{thm:main} is now proved. We derive the two consequences
stated in the introduction directly from its finite host. The first
uses tensor powers to amplify the homomorphism deficit and blow-ups
to compare injective copies at a positive limiting density. The second
keeps the edge density fixed while interpolating to an exact graphon
equality, and then samples a nonquasirandom sequence of finite graphs.

\subsection{Amplification and injective copies}

\begin{proof}[Proof of Corollary~\ref{cor:amplification}]
Take \(G\) from Theorem~\ref{thm:main} and put
\(N=|V(G)|\), \(p_0=p(G)\), \(t_0=t(H,G)\), and
\(\rho=t_0/p_0^{66}\). Simplicity and an edge give \(0<p_0<1\),
and mapping the two parts of \(H\) to the endpoints of an edge gives
\(t_0>0\). Thus \(0<\rho<1\).

Let \(G^{\times r}\) be the categorical power, whose adjacency is
coordinatewise. It is finite, simple and undirected, and
homomorphisms into it factor coordinatewise. Hence
\[
p(G^{\times r})=p_0^r,\qquad
t(H,G^{\times r})=t_0^r,\qquad
\frac{t(H,G^{\times r})}{p(G^{\times r})^{66}}=\rho^r\longrightarrow0,
\]
which proves (i). This is the standard tensor-power principle
\cite[Section~2]{conlonfoxsudakov2010}.

For (ii), use instead the balanced independent blow-ups \(F_s=G[s]\):
replace each vertex by an independent set of size \(s\), and each edge
by the complete bipartite graph between its two sets. Then \(n_s=Ns\),
\(p(F_s)=p_0\), and \(t(H,F_s)=t_0\), since each homomorphism into
\(G\) has exactly \(s^{35}\) lifts. For \(n_s\ge35\), set
\(q_s=\operatorname{Inj}(H,F_s)/(n_s)_{35}\). This is the homomorphism
probability for a uniform vertex map conditioned on injectivity, so
the collision union bound gives
\[
|q_s-t_0|\le1-\frac{(n_s)_{35}}{n_s^{35}}
\le\frac{\binom{35}{2}}{n_s}.
\]
Meanwhile \(d_s=p_0n_s/(n_s-1)\to p_0\). Therefore
\(q_s/d_s^{66}\to\rho<1\), proving (ii) with
\(p_*=p_0\) and \(\delta_*=(1-\rho)/2\).
\end{proof}

\subsection{Failure of the forcing conjecture}

\begin{proof}[Proof of Corollary~\ref{cor:forcing}]
Let \(G\) be the host from Theorem~\ref{thm:main}, put
\(N=|V(G)|\), and set \(p=p(G)\). Since \(G\) has an edge and is
simple,
\[
0<p\le\frac{N-1}{N}<1.
\]
Let \(W_0\) be its adjacency kernel on the uniform probability space
\(V(G)\). Independent vertex samples count all homomorphisms, so
\[
t(K_2,W_0)=p,\qquad t(H,W_0)=t(H,G)<p^{66}.
\]
Moreover, \(W_0\) is \(\{0,1\}\)-valued, and hence
\(\|W_0-p\|_2^2=p(1-p)>0\).

On the balanced sign space \(\{-1,1\}\), choose
\(0<\delta\le\min(p,1-p)\) and define the rank-one perturbation
\[
W_1(a,b)=p+\delta ab.
\]
It takes values in \([0,1]\) and has mean \(p\). For
\(A\subseteq E(H)\), let \(d_A(v)\) be the degree of \(v\) in the
edge set \(A\). Expanding the product defining \(t(H,W_1)\) and
averaging the independent signs gives
\[
t(H,W_1)=
\sum_{\substack{A\subseteq E(H)\\ d_A(v)\text{ even for every }v}}
p^{66-|A|}\delta^{|A|}>p^{66}.
\]
Indeed, the empty edge set contributes \(p^{66}\), every surviving
term is nonnegative, and the edge set of a cycle contributes a
strictly positive term. Such a cycle exists because a forest on
\(35\) vertices has at most \(34\) edges, whereas \(H\) has \(66\).

Lift the two kernels to the independent coordinates of
\(\Omega=V(G)\times\{-1,1\}\), with product probability, and set
\[
W_s\bigl((x,a),(y,b)\bigr)
=(1-s)W_0(x,y)+sW_1(a,b),\qquad 0\le s\le1.
\]
Each \(W_s\) is symmetric, takes values in \([0,1]\), and has mean
\(p\). The endpoint lifts preserve their \(H\)-densities. The two
centered summands depend on independent coordinate pairs and have
mean zero, so they are orthogonal in \(L^2(\Omega^2)\). Thus
\[
\|W_s-p\|_2^2=(1-s)^2p(1-p)+s^2\delta^2>0
\qquad(0\le s\le1).
\]
The function \(s\mapsto t(H,W_s)\) is a polynomial. Its endpoint
values lie strictly below and above \(p^{66}\), respectively, so
some \(s_*\in(0,1)\) satisfies \(t(H,W_{s_*})=p^{66}\). The norm
identity shows that \(W=W_{s_*}\) is nonconstant.

For completeness, this gives failure of asymptotic quasirandomness,
not just a graphon equality. On the finite atom space of \(W\), put
\[
r(x)=\E_z W(x,z),\qquad c(x,y)=\E_z W(x,z)W(y,z).
\]
Two applications of Cauchy--Schwarz give
\begin{align*}
t(C_4,W)&=\E_{x,y}c(x,y)^2\\
&\ge\bigl(\E_{x,y}c(x,y)\bigr)^2
=\bigl(\E_z r(z)^2\bigr)^2\ge p^4.
\end{align*}
Equality would force \(r(x)=p\) for every \(x\) and
\(c(x,y)=p^2\) for every pair \((x,y)\): every atom pair has positive
mass, including the diagonal pairs.
Then, for every \(x\),
\[
\E_z\bigl(W(x,z)-p\bigr)^2=c(x,x)-2pr(x)+p^2=0,
\]
contrary to nonconstancy. Hence \(t(C_4,W)>p^4\).

Use the independent sampling construction for this finite kernel,
as in Lemma~\ref{lem:finite-host} and
\cite[Section~2.6, Lemma~2.4]{lovaszszegedy2006}. For each fixed finite graph \(F\),
the collision count gives
\(\E t(F,G_n)=t(F,W)+O_F(n^{-1})\). The covariance of the indicators
for two vertex maps vanishes when their images are disjoint, and the
proportion of pairs with overlapping images is \(O_F(n^{-1})\).
Therefore \(\operatorname{Var}(t(F,G_n))=O_F(n^{-1})\). With tolerance
\(n^{-1/4}\), Chebyshev's inequality and a union bound for
\(F=K_2,H,C_4\) give total failure probability
\(O(n^{-1/2})<1\) for all sufficiently large \(n\). Choose one
successful realization for each such \(n\). The expectation estimate
then gives a deterministic sequence of finite simple graphs with
\[
p(G_n)\to p,\qquad t(H,G_n)\to p^{66},\qquad
t(C_4,G_n)\to t(C_4,W)>p^4.
\]
This sequence is not quasirandom of density \(p\), proving the claim.
\end{proof}

The graphon equality is exact, whereas the finite sequence matches the
edge and $H$-densities only in the limit. The argument supplies one
density $p\in(0,1)$ and does not assert failure of $p$-forcing at every
density.
